设平面图有 $n$ 个顶点、$m$ 条边、$f$ 个面（包含无界面），连通块数为 $c$。

\textbf{Euler 公式}

连通平面图满足
\[
  n-m+f=2.
\]
一般地，若有 $c$ 个连通块，则
\[
  n-m+f=1+c.
\]
所有面的边界长度之和为 $2m$；桥在同一个面的边界上出现两次，也应计数两次。

\textbf{简单平面图的边数上界}

若图连通且 $n\ge 3$，每个面的边界长度至少为 $3$，于是
\[
  2m\ge 3f,
  \qquad
  m\le 3n-6.
\]
极大简单平面图中每个面都是三角形，并且
\[
  m=3n-6,
  \qquad
  f=2n-4.
\]

若图的围长（最短环长度）为 $g$，则每个面的边界长度至少为 $g$，从而
\[
  m\le \frac{g}{g-2}(n-2).
\]
特别地，简单二分平面图无奇环，$g\ge4$，所以
\[
  m\le2n-4.
\]
森林没有有限围长，应直接使用 $m\le n-c$，不要套用围长公式。

以上边数上界对不连通图仍成立，只要满足相应的顶点数条件；分连通块计算通常能得到更紧的界。

\textbf{度数与染色}

由握手定理与 $m\le3n-6$，简单平面图的平均度严格小于 $6$，因此必有度数不超过 $5$ 的顶点。每个子图仍是平面图，所以平面图是 $5$-退化图，可按退化序贪心使用至多 $6$ 种颜色。

进一步有五色定理与四色定理：每个简单平面图分别一定可五染色、四染色。竞赛中若只需容易实现的保证，通常使用退化序得到六染色。

简单二分平面图的平均度严格小于 $4$，因此必有度数不超过 $3$ 的顶点。

\textbf{对偶图}

固定一个连通平面嵌入：原图的每个面对应对偶图的一个顶点；原图每条边连接其两侧的面，得到一条对偶边。桥对应对偶图中的自环，原图的重边可能对应对偶图的重边。

在连通平面图中，原图的一组边构成生成树，当且仅当这些边在对偶图中的补集构成对偶图的生成树。平面图中的环与对偶图中的割相对应，最小割问题有时可转为对偶图上的最短路。

\textbf{不可平面性的常用必要条件}

若简单图满足 $m>3n-6$，则一定不是平面图；若它还是二分图且 $m>2n-4$，也一定不是平面图。反之不成立，满足边数上界不代表图一定可平面。

Kuratowski 定理：一个有限图是平面图，当且仅当它不含 $K_5$ 或 $K_{3,3}$ 的细分作为子图。因此 $K_5$ 与 $K_{3,3}$ 是最常见的不可平面性证据。
